\chapter{Sauberes Kapitel}
\label{chap:sauber}

Dieses Kapitel enthält ausschließlich regelkonforme Formulierungen.
Es dient dazu, Fehlalarme des Linters sichtbar zu machen.

\section{Typografie und Zahlen}
\label{sec:sauber_typo}

Das betrifft z.\,B. die Stimme und d.\,h. auch das Tempo.
Der Anteil liegt bei 15\,\% der Beiträge, also bei $N = 40$ Personen.
Die Seiten S.~12--15 behandeln den Zeitraum 2020--2026.
Kapitel 5 beschreibt Item 3 und Hypothese 2 genauer.
Die Antworten liegen auf einer Skala von 1 bis 5.
Die Pause dauert 8\,s und das Tempo steigt um 5\,\%.
Es nahmen zwischen 5 und 40 Personen teil.
Die Studie umfasst drei Phasen und zwölf Items.
Der Begriff \enquote{Persona Collapse} stammt aus der Fachsprache.

\section{Sprache und Belege}
\label{sec:sauber_sprache}

Die Ergebnisse zeigen, dass die Adaption wirkt.
Der Unterschied ist signifikant (Wilcoxon-Test, $p = 0{,}03$).
\textcite{goldberg1990} beschreibt fünf breite Faktoren der Persönlichkeit.
Studien zeigen, dass Extraversion hörbar ist \parencite{goldberg1990}.
Ein \gls{a:KI}-Voice-Agent passt sich der Person an.
Der Nutzer bewertet den Agenten nach dem Gespräch.
Weitere Befunde liefern \textcite{ohnedoi2021} und \textcite{preprint2025}.

\begin{quote}
Wir haben das Gespräch als angenehm empfunden.
\end{quote}

\section{Abbildungen und Formeln}
\label{sec:sauber_abb}

Den Ablauf zeigt \cref{fig:sauber}.
Die Darstellung folgt einer bekannten Vorlage.
Sie wurde für die Arbeit angepasst.

\begin{figure}[htbp]
    \centering
    \rule{2cm}{1cm}
    \caption{Ablauf. Eigene Darstellung nach \textcite{goldberg1990}}
    \label{fig:sauber}
\end{figure}

Die Werte fasst \cref{tab:sauber} zusammen.
Sie nutzt den booktabs-Stil.
Sie enthält zwei Zeilen.

\begin{table}[htbp]
    \centering
    \caption{Werte}
    \label{tab:sauber}
    \begin{tabular}{lr}
        \toprule
        A & 1 \\
        \midrule
        B & 2 \\
        \bottomrule
    \end{tabular}
\end{table}

Der fusionierte Wert ergibt sich als gewichtetes Mittel.
\begin{equation}
    s = w_1 x_1 + w_2 x_2
\end{equation}
Dabei bezeichnet $s$ den fusionierten Wert, $w_1$ und $w_2$ die Gewichte und $x_1$ und $x_2$ die Schätzungen.
Ein Beispiel: Mit $w_1 = 1$, $w_2 = 0{,}5$, $x_1 = 0{,}4$ und $x_2 = 0{,}2$ ergibt sich $s = 0{,}5$.
